\chapter{Contexto y Estado del Arte}

% El título anterior era «Fundamentos de computación cuántica». Ya no vale: el capítulo
% cubre ahora también los modelos de aprendizaje automático. De paso desaparece el
% Overfull \hbox de 3,35 pt que obligaba a partir el título a mano.
\section{Contexto del problema}
\label{cap:fundamentos}

% =====================================================================================
% GUION. Capitulo de CONTEXTO: recoge conocimiento que YA EXISTE y se cita.
% Reestructurado en ago-2026 en dos bloques, 2.1 cuantico y 2.2 aprendizaje automatico.
%
% 🔴 LA FRONTERA CON EL CAPITULO 9, que es la regla que evita duplicar:
%
%      2.2  → QUE ES el metodo. Conocimiento de la literatura. Se CITA.
%      cap. 9 → POR QUE ese metodo aqui, como se configura, que se midio.
%               Decisiones propias. Se JUSTIFICA con nuestros datos.
%
%    Ejemplo: «que es la regularizacion L2» va en 2.2; «por que Ridge y no minimos
%    cuadrados, dada la colinealidad extrema que medimos» va en el 9. Si una frase
%    contiene una cifra nuestra, esta en el capitulo equivocado.
%
% 🔴 TONO, y es asimetrico a proposito. El tribunal es de INTELIGENCIA ARTIFICIAL:
%      2.1 → no saben cuantica. Claro, ameno, con intuicion antes que formalismo.
%            El tecnicismo entra, pero siempre despues de haberlo hecho comprensible.
%      2.2 → hablas su idioma. Denso y rapido en Ridge y Random Forest, que conocen
%            mejor que nadie. Detenerse SOLO en la red de grafos, que es lo que
%            probablemente no han visto aplicado a esto.
%
% ⚠️ Reglamento art. 11.4.a: «no es original el trabajo que contenga un exceso de
%    informacion de otras fuentes, AUNQUE ESTEN DEBIDAMENTE CITADAS». En un capitulo
%    de fundamentos es donde mas facil es caer. Defensa: incluir solo lo que los
%    capitulos 3, 6 y 9 van a usar de verdad, y que cada apartado apunte hacia adelante.
%
% ⚠️ Solo se numera hasta \subsection (secnumdepth = 2 en la clase book). Un
%    \subsubsection saldria SIN numero. Por eso el cuarto nivel va como parrafo con
%    entradilla en negrita, igual que en el resto de la memoria.
% =====================================================================================

% Parrafo de entrada del capitulo: que hace falta saber para leer el resto y por que
% este capitulo tiene dos mitades que parecen inconexas y no lo son.
% ⚠️ La norma prohibe dos encabezados consecutivos sin texto entre medias.


% =====================================================================================

% =====================================================================================
% ⚠️ DECISION DE REDACCION (sep-2026): la seccion declara su fuente UNA SOLA VEZ, en el
%    parrafo de entrada, y las citas puntuales a Nielsen y Chuang se retiran del cuerpo.
%    Motivo: la seccion entera es exposicion de un unico manual de referencia; repetir la
%    misma cita en cada parrafo no aporta trazabilidad y entorpece la lectura.
%    ⚠️ La atribucion del pie de la Figura 2.2 SI se conserva: es la fuente de la figura,
%       no una cita de teoria, y sin ella la figura quedaria sin origen.
%    ⚠️ Si en el futuro entra en 2.1 material de OTRA fuente, ese pasaje si lleva su cita.

Salvo indicación expresa, toda la teoría de esta sección está inspirada del manual de referencia de
la disciplina, \emph{Quantum Computation and Quantum Information} de
\citeA{nielsen2010quantum}.

\subsection{Del bit al qubit}
\label{sec:bit_qubit}
% LA PUERTA DE ENTRADA. Si el tribunal se pierde aqui, se pierde en todo el capitulo.
% Arrancar por lo que SI conocen: el bit, y de ahi la diferencia.
%
% Fuente: Nielsen y Chuang (2010), seccion 1.2, p. 13.
% ⚠️ VERIFICAR LA PAGINA del ultimo parrafo (el estado |+> y el ejemplo del atomo
%    iluminado): puede estar en la p. 14 y no en la 13. No citar p. 13 para eso sin
%    comprobarlo en el libro.
El bit es el concepto fundamental de la computación y la información clásicas. La
computación y la información cuánticas se construyen sobre un concepto análogo, el bit
cuántico o, abreviadamente, el \textbf{qubit}.

Igual que un bit clásico tiene un estado (o $0$, o $1$), un qubit
también tiene un estado. Dos estados posibles para un qubit son $\ket{0}$ y $\ket{1}$, que,
como cabe suponer, se corresponden con los estados $0$ y $1$ de un bit clásico. La notación
$\ket{\ }$ se denomina \textit{notación de Dirac} y es la
notación estándar para los estados en mecánica cuántica. La diferencia entre bits y qubits
es que un qubit puede encontrarse en un estado distinto de $\ket{0}$ o $\ket{1}$: también es
posible formar combinaciones lineales de estados, llamadas habitualmente
\textbf{superposiciones}:
\begin{equation}
\label{eq:superposicion}
\ket{\psi} = \alpha \ket{0} + \beta \ket{1}.
\end{equation}
Los números $\alpha$ y $\beta$ son números complejos.

Un bit se puede examinar para determinar si está en el estado $0$ o en el $1$. Los ordenadores
lo hacen constantemente, por ejemplo cuando recuperan el contenido de su memoria. Cabe recalcar
que un qubit no se pueda examinar para determinar su estado cuántico, es decir, los
valores de $\alpha$ y $\beta$. En su lugar, la mecánica cuántica nos dice que solo podemos
adquirir una información mucho más restringida sobre el estado cuántico. Al medir un qubit
obtenemos o bien el resultado $0$, con probabilidad $|\alpha|^2$, o bien el resultado $1$, con
probabilidad $|\beta|^2$. Naturalmente, $|\alpha|^2 + |\beta|^2 = 1$, puesto que las
probabilidades deben sumar uno.

Que un qubit pueda estar en un estado de superposición choca con nuestra comprensión de
«sentido común» del mundo físico que nos rodea. Un bit clásico es como una moneda: o cara,
o cruz. En monedas imperfectas puede haber estados intermedios, como quedarse en equilibrio
sobre el canto, pero en el caso ideal se pueden descartar. Un qubit, por el contrario, puede
existir en un continuo de estados entre $\ket{0}$ y $\ket{1}$, hasta que se lo observa.

Los qubits no son una idea
abstracta. Se han comprobado en muchos experimentos y existen distintas formas de
construirlos en el mundo real:
\begin{itemize}
    \item un qubit puede ser la polarización de un fotón, esto es, dos formas distintas de
          vibrar la luz;
    \item puede ser la orientación del espín de un núcleo dentro de un campo magnético;
    \item puede ser dos estados de un electrón en un átomo: el estado «fundamental»
          ($\ket{0}$) y el estado «excitado» ($\ket{1}$).
\end{itemize}

En este último ejemplo, el átomo funciona como un pequeño sistema de dos niveles. Si se le
ilumina con luz de la energía adecuada, se puede hacer que el electrón pase de $\ket{0}$ a
$\ket{1}$, o de $\ket{1}$ a $\ket{0}$. Hasta aquí, todo es clásico: dos estados, uno u otro.

Pero lo interesante llega después. Si no se le da suficiente tiempo de luz, el electrón no
llega completamente a $\ket{1}$, sino que se queda en un estado intermedio, una superposición
entre $\ket{0}$ y $\ket{1}$. Ese estado intermedio es lo que se denomina $\ket{+}$.

Un qubit se puede definir geométricamente sobre una esfera tridimensional. Esta esfera se denomina habitualmente
\textbf{esfera de Bloch}, y nos permite visualizar el estado de un sol qubit. Se adjunta la Figura~\ref{fig:bloch} para tener un idea visual del concepto del qubit.

\begin{figure}[htbp]
    \centering
    % =====================================================================================
% Esfera de Bloch — figura propia en TikZ.
%
% Se incluye desde capitulos/02_fundamentos.tex con \input. Es vectorial, usa la
% tipografia del documento y es de ELABORACION PROPIA: no arrastra ninguna licencia.
%
% GEOMETRIA. El plano ecuatorial se dibuja como una elipse alineada con los ejes de
% la pagina, de semiejes \R y \K*\R. Un punto del ecuador se parametriza con el angulo
% de pantalla t:            (\R cos t , \K \R sin t)
% Los ejes se colocan sobre esa misma elipse, asi que todo queda consistente:
%     \Ty  = 0°    -> el eje y, hacia la derecha
%     \Tx  = 235°  -> el eje x, hacia abajo a la izquierda
% El angulo azimutal phi se mide DESDE el eje x hacia el eje y, luego t = \Tx + phi.
%
% PARAMETROS AJUSTABLES: \R (radio), \K (aplastamiento de la perspectiva),
% \Theta y \Phi (posicion del estado dibujado).
% =====================================================================================
% ⚠️ \shorthandoff{<>} ES NECESARIO. `babel` con la opcion `spanish` convierte < y >
%    en caracteres ACTIVOS (los usa para las comillas << >>), y eso rompe la sintaxis
%    de puntas de flecha de TikZ: «>={Stealth}» y «[->]» fallan con
%    «Argument of \language@active@arg> has an extra }».
%    Va dentro de un grupo, asi que los catcodes se restauran solos al cerrar.
\begingroup
\shorthandoff{<>}%
\begin{tikzpicture}[>={Stealth[length=2.2mm]}, line width=0.5pt]

  \pgfmathsetmacro{\R}{2.2}      % radio de la esfera
  \pgfmathsetmacro{\K}{0.30}     % achatamiento del ecuador en perspectiva
  \pgfmathsetmacro{\Tx}{235}     % angulo de pantalla del eje x
  \pgfmathsetmacro{\Theta}{40}   % angulo polar del estado dibujado
  \pgfmathsetmacro{\Phi}{55}     % angulo azimutal del estado dibujado

  % --- coordenadas del estado |psi> ------------------------------------------------
  \pgfmathsetmacro{\Tp}{\Tx+\Phi}                       % angulo de pantalla del azimut
  \pgfmathsetmacro{\Rxy}{\R*sin(\Theta)}                % radio de la proyeccion ecuatorial
  \pgfmathsetmacro{\Px}{\Rxy*cos(\Tp)}
  \pgfmathsetmacro{\Pyeq}{\K*\Rxy*sin(\Tp)}             % altura de la proyeccion
  \pgfmathsetmacro{\Py}{\Pyeq+\R*cos(\Theta)}           % altura del propio estado

  % =====================================================================================
  % EL BIT CLASICO, a la izquierda, para la comparativa.
  % ⚠️ NO TOCA NADA DE LA ESFERA: se dibuja en su propia columna, a la izquierda del
  %    origen. La esfera sigue centrada en (0,0) exactamente como estaba.
  %
  % Los dos puntos van al MISMO azul que los polos de la esfera, y alineados con ellos.
  % Es deliberado: asi se lee de un vistazo que el bit clasico tiene DOS estados y que
  % el qubit tiene esos mismos dos... y toda la superficie que hay entre medias.
  % =====================================================================================
  \pgfmathsetmacro{\Xbit}{-2.60*\R}     % columna del bit clasico
  % ⚠️ Los cuatro puntos van al MISMO radio (\Rpunto) y las dos parejas a la MISMA
  %    altura, y=+\R y y=-\R: asi el 0 queda alineado con |0> y el 1 con |1>.
  \pgfmathsetmacro{\Rpunto}{3.2}        % radio de los puntos, en pt

  \fill[rojoclasico] (\Xbit,\R) circle (\Rpunto pt);
  \node[right=5pt] at (\Xbit,\R) {$0$};
  \fill[rojoclasico] (\Xbit,-\R) circle (\Rpunto pt);
  \node[right=5pt] at (\Xbit,-\R) {$1$};

  % rotulos de cada mitad
  \node[font=\bfseries] at (\Xbit,{-1.62*\R}) {Bit clásico};
  \node[font=\bfseries] at (0,{-1.62*\R}) {Qubit};

  % --- esfera y ecuador -------------------------------------------------------------
  \draw (0,0) circle (\R);
  \draw[dashed] (0,0) ellipse ({\R} and {\K*\R});

  % --- ejes -------------------------------------------------------------------------
  % Las puntas de flecha MUEREN sobre la propia esfera, sin sobresalir; las etiquetas
  % van fuera. El eje x es mas corto a proposito: no apunta hacia el borde del circulo
  % sino hacia el punto del ECUADOR que le corresponde, que en perspectiva cae dentro.
  % eje z: atraviesa la esfera de polo a polo, con punta arriba
  \draw (0,-\R) -- (0,0);
  \draw[->] (0,0) -- (0,\R);
  % El eje z lleva SU PROPIA etiqueta en negro, como x e y. Las dos etiquetas azules
  % de arriba y abajo no nombran el eje: nombran los dos ESTADOS de la base.
  \node[right=2pt] at ({0.02*\R},{0.86*\R}) {$\hat{\mathbf{z}}$};
  % Solo el estado: el eje ya lo nombra la etiqueta negra de arriba. Centradas sobre
  % el polo, que es lo que marcan.
  \node[above=2pt, azulunir] at (0,\R) {$\ket{0}$};
  \node[below=2pt, azulunir] at (0,-\R) {$\ket{1}$};
  % eje y: hacia la derecha, hasta el borde
  \draw[->] (0,0) -- (\R,0);
  \node[right=2pt] at (\R,0) {$\hat{\mathbf{y}}$};
  % eje x: hasta su punto del ecuador, abajo a la izquierda
  \draw[->] (0,0) -- ({\R*cos(\Tx)},{\K*\R*sin(\Tx)});
  \node[below left, inner sep=1pt] at ({\R*cos(\Tx)},{\K*\R*sin(\Tx)})
        {$\hat{\mathbf{x}}$};

  % --- los dos estados de la base computacional --------------------------------------
  % |0> y |1> son EXACTAMENTE los polos de la esfera. Marcarlos es correcto y ayuda:
  % deja ver de un vistazo que los estados clasicos son solo dos puntos de todos los
  % posibles. En azul institucional, el mismo de los titulos de capitulo.
  \fill[azulunir] (0,\R) circle (\Rpunto pt);
  \fill[azulunir] (0,-\R) circle (\Rpunto pt);

  % --- proyeccion del estado sobre el plano ecuatorial -------------------------------
  \draw[dotted] (0,0) -- (\Px,\Pyeq);
  \draw[dotted] (\Px,\Pyeq) -- (\Px,\Py);

  % --- el estado --------------------------------------------------------------------
  \draw[->] (0,0) -- (\Px,\Py);
  \fill (\Px,\Py) circle (1.5pt);
  \node[above right, inner sep=2pt] at (\Px,\Py) {$\ket{\psi}$};

  % --- angulo theta: entre el eje z y el estado --------------------------------------
  % El arco NO va de 90 a 90-theta: theta es el angulo real en 3D, y sobre el papel la
  % direccion del estado esta en atan2(\Py,\Px). Hay que arcar hasta ESE angulo para que
  % el arco muera justo sobre la flecha.
  \pgfmathsetmacro{\Apsi}{atan2(\Py,\Px)}
  \pgfmathsetmacro{\Rt}{0.46*\R}
  \draw (0,\Rt) arc[start angle=90, end angle=\Apsi, radius=\Rt];
  \pgfmathsetmacro{\Amid}{0.5*(90+\Apsi)}
  \node at ({1.24*\Rt*cos(\Amid)},{1.24*\Rt*sin(\Amid)}) {$\theta$};

  % --- angulo phi: en el plano ecuatorial, del eje x a la proyeccion -----------------
  % Se dibuja sobre la MISMA elipse del ecuador, a radio reducido, para que se lea como
  % un angulo tumbado en el plano y no como un angulo en el plano del papel.
  \pgfmathsetmacro{\Rp}{0.42*\R}
  \draw plot[domain=\Tx:\Tp, samples=40, smooth]
        ({\Rp*cos(\x)},{\K*\Rp*sin(\x)});
  \pgfmathsetmacro{\Pmid}{\Tx+0.5*\Phi}
  \node at ({1.55*\Rp*cos(\Pmid)},{1.55*\K*\Rp*sin(\Pmid)}) {$\phi$};

\end{tikzpicture}%
\endgroup


    \caption{Un bit clásico frente a un qubit: dos estados posibles frente
             a estados infinitos sobre la esfera de Bloch.}
    \label{fig:bloch}

    {\footnotesize Fuente: Claude Code \cite{anthropic_claude}.}
\end{figure}
% PENDIENTE en esta misma subseccion, aun sin redactar:
%   - Definir |+> explicitamente, si se decide dar la formula: (|0> + |1>)/raiz(2).
%   - El entrelazamiento como segunda propiedad que da la ventaja computacional.
%   - Por que n qubits necesitan 2^n amplitudes, y por que eso es a la vez el argumento
%     que hace interesante el campo y el que limita el tamaño de nuestro conjunto de
%     datos (capitulo~\ref{cap:planteamiento}).
%   - La contrapartida: el estado cuantico es
%     fragil, y esa fragilidad es el problema que este trabajo intenta mitigar.

\subsection{Qubits múltiples}
\label{sec:multiples}
% Apartado BREVE y con poca matematica: su unico objetivo es que se vea el crecimiento
% 2^n antes de hablar de puertas, porque despues se usa constantemente.
% Fuente: Nielsen y Chuang (2010), seccion 1.2.
% ⚠️ VERIFICAR LA PAGINA antes de la entrega. Se cita por seccion.
Un solo qubit no basta para calcular nada interesante. Entonces, ¿cómo
se describe el estado conjunto de varios? La respuesta explica por sí sola por
qué un ordenador cuántico puede hacer cosas que uno clásico no.

Con \textbf{dos qubits} las combinaciones posibles de valores clásicos son cuatro: $00$, $01$,
$10$ y $11$. A cada una le corresponde un estado de la base, y el estado conjunto es una
superposición de los cuatro, con una amplitud compleja por cada uno:
\begin{equation}
\label{eq:dos_qubits}
\ket{\psi} = \alpha_{00}\ket{00} + \alpha_{01}\ket{01}
           + \alpha_{10}\ket{10} + \alpha_{11}\ket{11}
\
\end{equation}

Con \textbf{tres qubits} el razonamiento se repite. Las cadenas posibles pasan a ser ocho
(de $000$ a $111$), y el estado necesita ocho amplitudes:
\begin{equation}
\label{eq:tres_qubits}
\ket{\psi} = \alpha_{000}\ket{000} + \alpha_{001}\ket{001} + \cdots + \alpha_{111}\ket{111}
\;\equiv\;
\begin{pmatrix} \alpha_{000} & \alpha_{001} & \cdots & \alpha_{111}\end{pmatrix}^{\!\intercal}.
\end{equation}

Y así sucesivamente:
\begin{itemize}
    \item con \textbf{cuatro qubits}, los estados de la base van de $\ket{0000}$ a
          $\ket{1111}$ y hacen falta \textbf{$16$ amplitudes};
    \item con \textbf{cinco qubits}, van de $\ket{00000}$ a $\ket{11111}$ y hacen falta
          \textbf{$32$ amplitudes};
    \item \dots
\end{itemize}

El patrón está claro: \textbf{$n$ qubits requieren $2^n$ amplitudes complejas}. Con 50
qubits son más de mil billones de números, y describir ese estado en un ordenador clásico
es imposible. Ahí está la promesa de la computación cuántica.

\subsection{Puertas cuánticas y circuitos}
\label{sec:puertas}
% Fuente de los dos primeros parrafos: Nielsen y Chuang (2010), seccion 1.3.
% ⚠️ VERIFICAR LA PAGINA antes de la entrega. Se cita por seccion porque no esta
%    comprobada; no inventarla.
%
% ⚠️ Del original se ha DESCARTADO la promesa «including a circuit which teleports
%    qubits!»: anuncia un contenido que esta memoria no va a desarrollar, y prometer
%    algo que luego no aparece es de las cosas que un tribunal si nota.
Los cambios que experimenta un estado cuántico se pueden describir con el lenguaje de la
computación cuántica. Del mismo modo que un ordenador clásico se construye a partir de un
circuito eléctrico formado por cables y puertas lógicas, un ordenador cuántico se construye
a partir de un circuito cuántico formado por cables y puertas cuánticas elementales,
encargadas de transportar y manipular la información cuántica.

Los circuitos de un ordenador clásico constan de cables y puertas lógicas. Los cables
sirven para transportar la información a lo largo del circuito, mientras que las puertas
lógicas realizan manipulaciones sobre esa información y la convierten de una forma en otra.

% La forma vectorial ya se introdujo en el apartado 2.1.1; aqui solo se recupera.


% ⚠️ FORMATO: titulo arriba y fuente abajo (Instrucciones §1.3, p. 6).
\begin{table}[htbp]
    \centering
    \begin{tabular}{@{}llp{0.46\textwidth}@{}}
        \toprule
        \textbf{Puerta} & \textbf{Matriz} & \textbf{Efecto} \\
        \midrule
        $X$ & $\begin{pmatrix} 0 & 1 \\ 1 & 0\end{pmatrix}$
            & Intercambia $\ket{0}$ y $\ket{1}$. Es la negación cuántica. \\[1.1em]
        $Y$ & $\begin{pmatrix} 0 & -i \\ i & 0\end{pmatrix}$
            & Intercambia los estados e introduce además una fase. \\[1.1em]
        $Z$ & $\begin{pmatrix} 1 & 0 \\ 0 & -1\end{pmatrix}$
            & Deja $\ket{0}$ intacto y cambia el signo de $\ket{1}$. \\[1.1em]
        $H$ & $\dfrac{1}{\sqrt{2}}\begin{pmatrix} 1 & 1 \\ 1 & -1\end{pmatrix}$
            & Puerta de Hadamard: crea superposición a partir de un estado de la base. \\[1.1em]
        $S$ & $\begin{pmatrix} 1 & 0 \\ 0 & i\end{pmatrix}$
            & Aplica una fase de $\pi/2$ al estado $\ket{1}$. \\[1.1em]
        $T$ & $\begin{pmatrix} 1 & 0 \\ 0 & e^{i\pi/4}\end{pmatrix}$
            & Aplica una fase de $\pi/4$. \\[1.1em]
        $R_z(\theta)$ & $\begin{pmatrix} 1 & 0 \\ 0 & e^{i\theta}\end{pmatrix}$
            & Rotación de fase de ángulo arbitrario en torno al eje $z$. \\
        \bottomrule
    \end{tabular}

    \caption{Puertas cuánticas de un qubit más habituales.}
    \label{tab:puertas}

    {\footnotesize Fuente: elaboración propia.}
\end{table}

En la Figura~\ref{fig:hadamard} se puede observar geométricamente cómo afecta una puerta Hadamard a un qubit en la esfera de Bloch.

% ⚠️ FORMATO: titulo arriba y fuente abajo (Instrucciones §1.3, p. 6).
% ⚠️ FUENTE: es un redibujo de una figura de Nielsen y Chuang, no una invencion propia.
%    La forma academica correcta es «elaboracion propia a partir de…», no «elaboracion
%    propia» a secas. Verificar el numero de figura del libro y añadirlo si se quiere.
% Codigo de color: azul = el qubit · naranja = la transformacion.
\begin{figure}[htbp]
    \centering
    % =====================================================================================
% La puerta de Hadamard vista sobre la esfera de Bloch — figura propia en TikZ.
%
% Tres viñetas: el estado de partida (|0>+|1>)/raiz(2) sobre el eje x, la rotacion de
% 90 grados en torno a y que lo lleva al polo sur, y la rotacion de 180 grados en torno
% a x que lo sube al polo norte. Resultado: H|+> = |0>.
%
% CODIGO DE COLOR, deliberado y constante en toda la memoria:
%   azulunir   -> el QUBIT (vector de estado, punto y etiqueta)
%   acentofig  -> las TRANSFORMACIONES (giros, planos, flechas de operacion)
%
% ⚠️ \shorthandoff{<>} es obligatorio: `babel` en español hace activos < y >, y eso
%    rompe la sintaxis de puntas de flecha de TikZ.
% =====================================================================================
\begingroup
\shorthandoff{<>}%
\begin{tikzpicture}[>={Stealth[length=2mm]}, line width=0.5pt, scale=0.92]

  \pgfmathsetmacro{\R}{1.45}    % radio de cada esfera
  \pgfmathsetmacro{\K}{0.30}    % achatamiento del ecuador; ver nota al pie del fichero
  \pgfmathsetmacro{\Tx}{235}    % angulo de pantalla del eje x
  \pgfmathsetmacro{\D}{4.9}     % separacion entre viñetas

  % --- una esfera con sus tres ejes, reutilizable ------------------------------------
  \newcommand{\esferabase}{%
    \draw (0,0) circle (\R);
    \draw[dashed] (0,0) ellipse ({\R} and {\K*\R});
    \draw[->] (0,0) -- (0,\R);
    \node[right=1pt] at ({0.02*\R},{0.78*\R}) {\scriptsize $\hat{\mathbf{z}}$};
    \draw[->] (0,0) -- (\R,0) node[right=-1pt] {\scriptsize $\hat{\mathbf{y}}$};
    \draw[->] (0,0) -- ({\R*cos(\Tx)},{\K*\R*sin(\Tx)})
          node[below left=-3pt] {\scriptsize $\hat{\mathbf{x}}$};%
  }

  % =====================================================================================
  % 1) Estado de partida: (|0>+|1>)/raiz(2), sobre el eje x
  % =====================================================================================
  \begin{scope}[xshift=0cm]
    \esferabase
    \draw[azulunir, line width=1.1pt]
          (0,0) -- ({\R*cos(\Tx)},{\K*\R*sin(\Tx)});
    \fill[azulunir] ({\R*cos(\Tx)},{\K*\R*sin(\Tx)}) circle (2pt);
    \node[azulunir, anchor=east] at ({-1.06*\R},{-0.40*\R})
         {\small $(\ket{0}+\ket{1})/\sqrt{2}$};
    % giro de 90 grados en torno al eje y
    \draw[acentofig, ->, line width=0.7pt]
          ({0.30*\R},{-0.038*\R}) arc[start angle=-170, end angle=110,
          x radius={0.14*\R}, y radius={0.22*\R}];
  \end{scope}

  \draw[->, line width=0.8pt] ({0.5*\D-0.45},0) -- ({0.5*\D+0.45},0);

  % =====================================================================================
  % 2) Tras el giro: |1> en el polo sur. La rotacion de 180 grados en torno a x se
  %    representa con el plano y la flecha doble.
  % =====================================================================================
  \begin{scope}[xshift=\D cm]
    % el plano va PRIMERO: asi la esfera y el estado quedan por encima y se ven
    \draw[acentofig, fill=acentofig, fill opacity=0.06, line width=0.6pt]
          ({-1.78*\R},{-0.33*\R}) -- ({-0.90*\R},{0.31*\R})
       -- ({ 1.78*\R},{ 0.31*\R}) -- ({ 0.90*\R},{-0.33*\R}) -- cycle;
    \esferabase
    \draw[azulunir, line width=1.1pt] (0,0) -- (0,-\R);
    \fill[azulunir] (0,-\R) circle (2pt);
    \node[azulunir, anchor=north] at (0,{-1.08*\R}) {\small $\ket{1}$};
    % flecha doble: el estado atraviesa el plano de arriba abajo
    \draw[acentofig, <->, line width=0.9pt]
          ({0.58*\R},{-0.52*\R}) -- ({0.58*\R},{0.52*\R});
  \end{scope}

  \draw[->, line width=0.8pt] ({1.5*\D-0.45},0) -- ({1.5*\D+0.45},0);

  % =====================================================================================
  % 3) Resultado: |0> en el polo norte
  % =====================================================================================
  \pgfmathsetmacro{\DD}{2*\D}
    \begin{scope}[xshift=\DD cm]
    \esferabase
    \draw[azulunir, line width=1.1pt] (0,0) -- (0,\R);
    \fill[azulunir] (0,\R) circle (2pt);
    \node[azulunir, anchor=south east] at ({-0.04*\R},{1.06*\R}) {\small $\ket{0}$};
  \end{scope}

\end{tikzpicture}%
\endgroup

% =====================================================================================
% NOTA SOBRE \K, EL ACHATAMIENTO
%
% Determina desde que altura se mira la esfera, y con ello el angulo VISUAL entre los
% ejes x e y. Con el eje y horizontal, ese angulo no puede llegar a 90 grados: el eje x
% ha de apoyarse sobre la elipse del ecuador —porque el ecuador ES el plano xy— y la
% unica forma de que forme 90 grados con y seria que apuntase recto hacia abajo, encima
% del eje z. Lo medido:
%
%     K = 0,30  ->  157 grados   punta de x a 0,62·R del centro
%     K = 0,55  ->  142 grados   punta de x a 0,73·R del centro   <- el elegido
%     K = 0,75  ->  133 grados   punta de x a 0,84·R del centro
%     K = 1,00  ->  125 grados   punta de x a 1,00·R  (el ecuador se confunde con
%                                el contorno y la esfera deja de parecer una esfera)
%
% Se probo 0,55 y se descarto: el usuario prefiere la perspectiva original de 0,30.
% La tabla se deja aqui por si en algun momento se quiere volver a mover.
% =====================================================================================


    \caption{Acción de la puerta de Hadamard sobre la esfera de Bloch.}
    \label{fig:hadamard}

    {\footnotesize Fuente: Claude Code \cite{anthropic_claude}, a partir de
    \citeNP{nielsen2010quantum}.}
\end{figure}

\paragraph{El circuito.} Un circuito cuántico es una secuencia de puertas aplicadas sobre
un conjunto de qubits, seguida de una medición. Se representa con una línea horizontal por
qubit (los «cables» del símil clásico) y las puertas colocadas sobre ellas en el orden en
que se ejecutan. Al circuito lo describen dos magnitudes: la \textbf{anchura},
que es el número de qubits, y la \textbf{profundidad}, que es el número de capas de puertas
que deben ejecutarse una tras otra.

La Figura~\ref{fig:circuito_ghz} recoge tres ejemplos, los tres de anchura $5$, que sirven
para ver hasta qué punto circuitos del mismo tamaño pueden ser muy distintos por dentro.

% ⚠️ FORMATO: titulo arriba y fuente abajo (Instrucciones §1.3, p. 6).
% Figura generada por el propio proyecto: TFM-Quantum/figures/00_anatomia/circuitos/.
% Se copia el PDF, no el PNG, porque es vectorial y no pixela.
%
% ⚠️ La imagen trae su PROPIO titulo incrustado («GHZ — Estado GHZ — el entrelazamiento
%    en su forma mas pura»), asi que se duplica con el \caption. Lo suyo es regenerarla
%    sin titulo desde src/viz/circuitos, o recortarlo. Se deja asi de momento.
\begin{figure}[htbp]
    \centering

    % Composicion compacta: GHZ y HEA en una fila, y la QFT —mas larga— debajo,
    % ocupando el ancho completo. Sin paquetes nuevos: dos minipages y un salto.
    % ⚠️ Las tres imagenes traen SU PROPIO titulo incrustado, que chocaba con el
    %    \caption y quedaba descolgado a la izquierda. Se recortan los 20 bp
    %    superiores de cada PDF —medido: el rotulo ocupa el 6,3 % del alto—.
    %    Lo limpio seria regenerarlas sin titulo desde TFM-Quantum/src/viz/circuitos.
    \includegraphics[trim=0 0 0 20bp, clip, width=0.47\textwidth]{cap02/circuito_ghz}\hfill
    \includegraphics[trim=0 0 0 20bp, clip, width=0.47\textwidth]{cap02/circuito_hea}

    \vspace{0.4em}
    \includegraphics[trim=0 0 0 20bp, clip, width=0.88\textwidth]{cap02/circuito_qft}

    {\footnotesize Fuente: elaboración propia.}
    \caption{Tres ejemplos de circuito: preparación de un estado GHZ, un ansatz
             \textit{hardware-efficient} y una transformada cuántica de Fourier.}
    \label{fig:circuito_ghz}
\end{figure}

\paragraph{La base nativa.} Un procesador real no implementa todas las puertas de la
Tabla~\ref{tab:puertas}, sino un conjunto reducido a partir del cual puede construirse
cualquier otra. Antes de ejecutarse, todo circuito se \textbf{traduce} a ese conjunto
mediante un proceso llamado \textit{transpilación}. En los procesadores IBM Heron la base
nativa es $\{cz, id, rz, sx, x\}$.

Esa traducción evidencia más problema que aborda este trabajo. Un
circuito teóricamente tiene diez puertas puede convertirse en cientos al expresarse en
la base nativa, y cada puerta ejecutada es una oportunidad más de introducir
error. Cuanto más largo es el circuito traducido, menos señal sobrevive hasta la medición.

% =====================================================================================
% PENDIENTE en esta subseccion:
%   - FIGURA de un circuito sencillo (dos qubits, H + CNOT + medicion). Seria muy util
%     para un tribunal que no sabe de cuantica: ahora mismo el circuito se describe con
%     palabras pero no se ve. Se puede dibujar en TikZ, como la esfera de Bloch, o con
%     el paquete `quantikz` si se decide añadirlo al preambulo.
%   - Verificar la pagina exacta de la cita de Nielsen y Chuang.
% =====================================================================================

\subsection{El circuito como grafo acíclico dirigido}
\label{sec:dag}
% ⚠️ ESTA SUBSECCION ERA LA BISAGRA hacia el 2.2: cerraba el bloque cuantico con la idea
%    de que, si un circuito es un grafo, predecirlo es regresion sobre grafos. Al subirla
%    aqui quedan DOS apartados por debajo (hardware y utilidad), asi que la transicion ya
%    no cae junto al 2.2. Hay que reponerla al final de «Para que sirve un ordenador
%    cuantico», o el 2.2 arrancara en seco.

Un circuito se puede ver como un \textbf{grafo acíclico dirigido} (DAG, \textit{directed
acyclic graph}): cada puerta es un \textbf{nodo}, y hay una \textbf{arista} de una puerta a
otra. Es acíclico porque los estados cuánticos evolucionan hacia adelante en el tiempo.

\begin{figure}[htbp]
    \centering
    \includegraphics[width=0.62\textwidth]{cap02/dag_grafo}

    \caption{El grafo acíclico dirigido de un circuito de cinco qubits: cada puerta es un
             nodo y cada arista.}
    \label{fig:dag}

    {\footnotesize Fuente: elaboración propia.}
\end{figure}


Leído así, la profundidad del circuito deja de ser un recuento y pasa a ser el
\textbf{camino más largo} del grafo. Es la magnitud que determina cuánto tiempo permanece vivo el estado cuántico y,
por tanto, cuánto ruido acumula.

% ⚠️ FORMATO: titulo arriba y fuente abajo (Instrucciones §1.3, p. 6).
% La figura la genera figuras/cap02/gen_circuito_dag.py, que la produce con el propio
% Qiskit: el grafo NO es una ilustracion inventada, es la representacion interna que el
% transpilador manipula de verdad. El coloreado de los nodos (verde entradas, azul
% operaciones, rojo salidas) lo pone `dag_drawer` por si mismo.
%
% ⚠️ CAMBIO (sep-2026): se retiro el circuito que acompañaba al grafo por la izquierda
%    (figuras/cap02/dag_circuito.pdf). El PDF se conserva en la carpeta por si hiciera
%    falta reponerlo; para ello, envolver ambos \includegraphics en dos minipage de
%    0.44 y 0.52 \textwidth. El pie de la figura hay que devolverlo entonces a
%    «Un circuito y su grafo acíclico dirigido: ...».


\textbf{Y aquí es donde la física se convierte en un problema de aprendizaje automático:} si
un circuito es un grafo, predecir su comportamiento es un problema de regresión sobre
grafos.

% =====================================================================================
% =====================================================================================
% 2.1.6 — EL HARDWARE ACTUAL
% Sustituye al bloque de comentarios que habia bajo \subsection{El hardware actual}.
% Enfoque acordado: divulgativo («a modo curiosidad»), sin tecnicismos, y quedandose
% solo con los datos que el resto de la memoria necesita: la topologia limitada, la
% recalibracion diaria y el coste de simular. Las cifras van fechadas en 2026 y con cita.
% =====================================================================================

\subsection{El hardware actual}
\label{sec:hardware}

Conviene hacerse una idea de qué son físicamente las máquinas de las que habla esta
memoria, aunque solo sea para situar con una idea general el problema. Un procesador cuántico
superconductor como los de IBM es un chip del tamaño de una moneda que vive en el fondo de
un criostato: un cilindro de cobre dorado, de casi dos metros, cuya única función es
mantener el chip a unas centésimas de grado por encima del cero absoluto. Los qubits no se
programan escribiendo en una memoria, sino disparándoles pulsos de microondas por unos
cables coaxiales; cada puerta del circuito es un pulso. Y la medición no devuelve un
resultado, sino una estadística.

% Parrafo nuevo: explica QUE significa «repetir miles de veces y contar». Es el que
% justifica que el objetivo del TFM sea un valor esperado y no una salida determinista.
Al medir un qubit no
se lee el estado en el que estaba: el estado \textbf{se destruye} y en su lugar aparece un
cero o un uno, elegido al azar con la probabilidad que marcaba ese estado. Una única
ejecución, por tanto, no informa de nada; devuelve una cadena de ceros y unos que podría
haber salido de otra máquina completamente distinta. Lo que se hace en la práctica es
ejecutar el mismo circuito miles de veces seguidas (cada repetición se llama
\emph{shot}) y anotar cuántas veces sale cada resultado. De ese recuento se obtiene el
promedio que interesaba, y el circuito se ha ejecutado entero, desde cero, en cada una de
esas miles de repeticiones. La consecuencia de todo esto es que la precisión es un factor que se paga caro, ya que para reducir el error
en los resultados es necesario ejecutar más veces los circuitos, lo que aumenta coste de uso del procesador.

Merece de interés mencionar dos características del hardware que lo hacen distinto de un ordenador clásico:

\textbf{Primera: los qubits no se hablan todos con todos.} En un chip clásico cualquier
dato puede llegar a cualquier registro; en uno cuántico cada qubit está conectado solo a
dos o tres vecinos. Si un algoritmo necesita que interactúen dos qubits que no lo son, la
única salida es ir acercando el estado de uno hasta el otro con puertas intermedias, y cada
puerta añadida es ruido añadido.

\textbf{Segunda: el chip no es el mismo dos días seguidos.} Estas máquinas se recalibran a
diario, y cada calibración publica un boletín con la tasa de error de cada puerta y los
tiempos de vida de cada qubit. Esos números cambian y un qubit que ayer era fiable hoy puede no serlo, y hay puertas que se estropean y no se reparan.
Esta es la razón de que el conjunto de datos del Capítulo~\ref{cap:planteamiento}
contenga la calibración real de un intervalo de días distintos en lugar de una única
fotografía del chip.

% ⚠️ FORMATO: titulo arriba y fuente abajo (Instrucciones §1.3, p. 6).
% Tabla simplificada a proposito: solo el catalogo de chips. El comentario sobre el
% cambio de estrategia (subir a 1121 y bajar a 156) se retiro por peticion expresa.
\begin{table}[htbp]
    \centering
    \begin{tabular}{@{}lrr@{}}
        \toprule
        \textbf{Familia} & \textbf{Año} & \textbf{Qubits} \\
        \midrule
        Falcon    & 2019 & 27 \\
        Eagle     & 2021 & 127 \\
        Osprey    & 2022 & 433 \\
        Condor    & 2023 & 1\,121 \\
        Heron     & 2023 & 133 / 156 \\
        Nighthawk & 2025 & 120 \\
        \bottomrule
    \end{tabular}

    \caption{Familias de procesadores cuánticos de IBM y su número de qubits.}
    \label{tab:familias}

    {\footnotesize Fuente: elaboración propia a partir de \citeNP{ibm2026hardware}.}
\end{table}

% Bloque nuevo: la equivalencia chip cuantico <-> ordenador clasico.
% Cifras calculadas como 2^n amplitudes x 16 bytes (complejo de doble precision).
% El 17,2 GB de n=15 es dato del propio proyecto: con ruido el coste es 4^n = 2^30.

Aunque $156$ qubits parecen poca cosa habría que repasar \textbf{cuánta memoria haría falta en un ordenador normal
para imitar el chip}. Y la respuesta es que describir el estado de $n$ qubits exige guardar
$2^n$ números, de modo que \textbf{cada qubit que se añade duplica} el coste. 
\begin{itemize}
    \item Con $30$ qubits hacen falta unos $17$ GB, la memoria de un portátil actual bueno.
    \item  Con $40$ qubits, unos $17$ TB, que ya es un servidor grande.
    \item Con $50$ qubits, unos $18$ PB: un centro de datos entero dedicado a imitar un chip del tamaño de una moneda. 
    \item Y con los $156$ qubits de un Heron sale un número de $48$ cifras que ninguna máquina actual ni futura podrá almacenar.
\end{itemize}

Estas ejecuciones con ruido de chips cuánticos se pueden simular con ordenadores clásicos mediante paquetes de 
simuladores, como Qiskit Aer. Debido al gran cómputo que exigen una ejecución cuántica, los 
circuitos simulados en este trabajo llegan hasta $15$ qubits.

Con todo esto hemos visto que el chip real maneja $156$ qubits sin problema y el ordenador clásico no es capaz
ni de acercarse. Por tanto, no se puede calcular de antemano cuánto va a degradar
el ruido un resultado, porque ese cálculo es precisamente lo que resulta imposible. Lo
que sí puede intentarse es \textbf{aprenderlo} a partir de ejemplos.



% =====================================================================================
\section{Estado del arte}
\label{cap:sota}
% =====================================================================================
% GUION. Reescrito en sep-2026 por indicacion del tutor: «algo escaso en numero de
% paginas, desarrollar mas las investigaciones previas».
%
% 🔴 SIN SUBAPARTADOS, por decision del autor. Antes eran cinco (§2.2.1 a §2.2.5) y se
%    fundieron en una sola seccion continua. NINGUNA de sus cinco etiquetas se citaba
%    desde fuera, asi que la fusion no rompio nada. La que SI se cita desde los capitulos
%    1, 3 y 7 es `cap:sota`, y se conserva.
%
% 🔴 EL HILO ES «CUANDO ACTUA CADA FAMILIA» respecto a la ejecucion del circuito, porque
%    es lo que separa a este trabajo del resto. No reordenar por tipo de modelo ni por
%    año: se perderia el argumento.
%
% 🔴 RESTRICCION DURA DE EXTENSION. La norma acota «contexto + estado del arte» a 10-15
%    paginas EN TOTAL. Con §2.1 en 7 paginas, esta seccion mas §2.3 pueden ocupar 7 como
%    mucho. MEDIR antes de añadir.
%
% ⚠️ LOS PREPRINTS VAN DECLARADOS COMO TALES. QuEst, SQuASH y Placidi no tienen actas
%    confirmadas —comprobado en dblp, ver doc/SoTA_ampliada/fichas—. Omitirlo seria
%    presentar como consolidado lo que no lo esta.
%
% Fuentes: doc/SoTA/comparative_analysis.md y doc/SoTA_ampliada/fichas/paper_{1..5}.md.
%
% ⚠️ Reglamento art. 11.4.a: «no es original el trabajo que contenga un exceso de
%    informacion de otras fuentes, AUNQUE ESTEN DEBIDAMENTE CITADAS». Defensa: de cada
%    trabajo se dice QUE DECIDE en los capitulos 4 a 6 de esta memoria.
%
% ⚠️ Excluidos por no citables: Wang/Zhengzhou, Stirbu, Al Farib.
% =====================================================================================

Este apartado realiza un recorrido de lo que ya existe para corregir el efecto del ruido, ordenado por
\textbf{cuándo} actúa cada familia de métodos respecto a la ejecución del circuito. Se ha
escogido este criterio porque es el que separa
a este trabajo de todo lo demás.

\subsubsection*{La mitigación clásica y su precio}

La corrección de errores completa (QEC) exige cientos de qubits físicos por cada qubit lógico,
inviable para las máquinas actuales \cite{preskill2018nisq}. La mitigación de errores (QEM) es
la alternativa: no elimina el ruido, sino que estima el resultado que se habría obtenido sin él
\cite{cai2023quantum}.

Las técnicas de referencia son la \textbf{extrapolación a ruido cero} (ZNE), que trata de estimar cuál habría sido el resultado de un circuito si el ruido del hardware no estuviera presente \cite{temme2017error}; la \textbf{cancelación probabilística de errores} (PEC), que utiliza métodos estadísticos para compensar los errores introducidos durante la ejecución del circuito, aunque requiere un mayor número de ejecuciones \cite{cai2023quantum}; y la \textbf{mitigación de lectura} (M3), centrada en corregir los errores que aparecen al medir el estado final de los qubits \cite{nation2023suppressing}. 


Los tres planteamientos comparten los dos rasgos que este trabajo toma como punto de partida:
todos \textbf{consumen ejecuciones adicionales} del procesador, que es el recurso más escaso, y
todos actúan \textbf{después} de ejecutar.

\subsubsection*{Sustitución de las ejecuciones por un modelo}

La primera forma de esquivar ese precio es reemplazar las ejecuciones extra por un modelo
entrenado. \citeA{liao2024machine} demostraron que una red puede aprender la corrección
directamente a partir del valor ruidoso medido, y lo validaron sobre hardware real de IBM hasta
un centenar de qubits. De su trabajo sale un dato que esta memoria aprovecha directamente:
entre los modelos tabulares que comparan, \textbf{Random Forest es la mejor referencia}, y por
eso es el que entra en la comparativa del Capítulo~\ref{cap:comparativa}.

Sobre esa misma idea, pero llevándola al grafo, GTraQEM \cite{bao2025gtraqem} representa el
circuito como un DAG (\textit{Directed Acyclic Graph}) y lo procesa con un \textit{graph
transformer} apoyado en un \textbf{nodo virtual} conectado a todos los demás, que recoge el
contexto global. Dos decisiones suyas se reutilizan aquí y una se descarta. Se reutiliza el nodo
virtual, por las razones que desarrolla el apartado~\ref{sec:gnn}. Se reutiliza también su
formulación del objetivo como \textbf{residuo} (predicen la diferencia entre el valor ruidoso y
el ideal, no el valor absoluto, porque esa diferencia es numéricamente pequeña y estabiliza el
entrenamiento), que es el mismo razonamiento que conduce al factor de supervivencia del
apartado~\ref{sec:factor}. Y se descarta su arquitectura: un \textit{transformer} sin paso de
mensajes, cuyo coste de atención era superior a la máquina disponible. QEMFormer \cite{bao2025qemformer} se encuentra entre los métodos más avanzados y combina información general del circuito con detalles sobre su estructura.

\textbf{Aquí aparece la diferencia que define esta memoria}: los tres toman el \textbf{valor
ruidoso medido como entrada}, de modo que hay que haber ejecutado el circuito para poder
corregirlo. \citeA{placidi2026deep} (preprint, sin actas
confirmadas) miden cuánto pesa esa entrada, y concluyen que es la que más
contribuye al rendimiento: sin ella, la precisión del modelo cae de forma considerable.

\subsubsection*{Técnicas pre-ejecución}

Existe una familia distinta que decide \textbf{antes} de gastar la máquina. La herramienta
oficial de IBM, mapomatic \cite{nation2023suppressing}, analiza distintas formas de asignar un mismo circuito al hardware disponible y selecciona la que previsiblemente producirá menos errores; es una
heurística analítica, sin aprendizaje. \citeA{hartnett2024learning} dan un paso más implementando técnicas de aprendizaje
automático con datos medidos en hardware real.

Una nueva linea de investigación sí predice magnitudes sobre el grafo del circuito antes de ejecutarlo, y es
lo más cercano a esta propuesta. Se destacan alguno de estos trabajos.

\textbf{QuEst} \cite{wang2025quest} (preprint, sin actas confirmadas) estima la fiabilidad
esperada de un circuito antes de ejecutarlo en la QPU. Construye el DAG, describe cada nodo con el
tipo de puerta, el qubit sobre el que actúa y la telemetría de ese qubit, y lo procesa con un
\textit{graph transformer} cuya atención se restringe a los vecinos del grafo. Reportan un
coeficiente de determinación de $0{,}991$ en simulador ruidoso y por encima de $0{,}95$ en cinco
procesadores reales de IBM, con una predicción unas doscientas veces más rápida que simular.
Sin embargo, los propios autores declaran que frente a una red tabular sobre variables
agregadas, la ventaja media es de $0{,}014$ de error cuadrático medio, pero \textbf{en tres de
las cinco máquinas reales la diferencia baja de $0{,}0034$ y en una de ellas gana la tabular}.
Casi toda la ventaja se concentra en los circuitos estructurados.

\textbf{Liu y colaboradores} \cite{liu2026output}, ya publicado en revista, predicen
directamente el valor esperado que producirá el circuito, con un coeficiente de determinación
entre $0{,}90$ y $0{,}998$ según el tamaño y la profundidad. La aportación más relevante es la siguiente: con el \emph{mismo} modelo y la \emph{misma}
arquitectura, cambiar qué magnitud se predice les mejora el acierto en un $36$\,\%. Es una
evidencia clara de que \textbf{reformular el objetivo puede cambiar de modelo}, que es
exactamente lo que mide el apartado~\ref{sec:candidatas} de esta memoria.

\textbf{Tudisco y colaboradores} \cite{tudisco2025graph}, en actas de IEEE QCE 2025, resuelven
otro problema (en qué máquina conviene ejecutar) pero su análisis previo es lo que interesa:
antes de proponer su red, \textbf{demuestran que el número de qubits, la profundidad y el
recuento de puertas no bastan} para separar los casos. De ahí concluyen que hay que meter el circuito entero en el modelo en vez de
resumirlo a mano. Es el mismo argumento que sostiene la comparativa de esta memoria.

\textbf{Aktar y colaboradores} \cite{aktar2024graph}, también en actas de IEEE QCE, predicen la
expresibilidad de circuitos parametrizados con una red de grafos adaptada de QuEst. Aportan dos
medidas que este trabajo utiliza. La primera es la evidencia en la \textbf{extrapolación en tamaño}:
un modelo entrenado hasta tres qubits se dispara a un error de $0{,}22$ al predecir seis, y uno
entrenado hasta ocho se queda en $0{,}05$ al predecir diez. La segunda conclusión es que \textbf{La variedad de condiciones de ruido durante el entrenamiento es más importante que entrenar con el perfil de ruido exacto de la máquina objetivo.}. Al evaluar sobre un dispositivo cuántico real, el modelo entrenado con tres perfiles de ruido diferentes casi duplica la tasa de acierto del modelo entrenado exclusivamente con el perfil de ruido de esa misma máquina. Este resultado aporta evidencia empírica a favor de entrenar con los cuarenta y dos días de calibración disponibles, ya que la variabilidad en las condiciones de ruido parece favorecer una mejor capacidad de generalización. De ellos sale además el precedente publicado de
la \textbf{pérdida de Huber} que usa el apartado~\ref{sec:gnn}.

\textbf{SQuASH} \cite{martyniuk2025benchmarking} (preprint) es el \textbf{precedente directo de
esta comparativa}: enfrenta una red de grafos a un Random Forest sobre los mismos circuitos y el
mismo protocolo, y la red gana en las tres familias que prueban, con coeficientes de $0{,}945$
frente a $0{,}705$ en el caso más favorable.

Ninguno de los cinco predice la degradación de un observable, y ninguno modela el paso del
tiempo.

\subsubsection*{El hardware cambia con el tiempo}

\citeA{hirasaki2023detection} midieron que las tasas de error no se degradan de forma
progresiva, sino \textbf{a saltos bruscos} que aparecen en cualquier momento y duran
minutos o decenas de minutos, y que la recalibración diaria de IBM no basta para capturarlos.
\citeA{huo2026anchor}, desde la perspectiva de sistemas, lo confirman por el otro extremo: el
mismo programa ejecutado en momentos distintos devuelve resultados inconsistentes, y proponen
una técnica que reduce esa variabilidad. Uno aporta la evidencia física del fenómeno y el otro
su impacto medible sobre el usuario final. Entre los dos justifican que el conjunto de datos de
esta memoria tenga un \textbf{eje temporal} con días de calibración reales.

% =====================================================================================
\section{Conclusiones}
\label{sec:gap}
% =====================================================================================
% 🔴 SECCION NUEVA (sep-2026, decision del autor). Recoge lo que antes era §2.2.5 «El
%    hueco en la literatura»: el hueco ES la conclusion del estado del arte, asi que la
%    seccion nace con contenido propio en vez de repetir lo anterior.
%    ⚠️ La etiqueta `sec:gap` se conserva desde alli. No renombrarla.
% =====================================================================================

% ⚠️ FORMATO: la plantilla oficial pone el \caption DEBAJO. Se sigue la plantilla.
% ⚠️ Va anclada con [t] y en \small A PROPOSITO: con [htbp] y cuerpo normal no cabia al
%    pie de su pagina y arrastraba la lista de debajo.
% Las cuatro columnas son los cuatro diferenciadores enunciados debajo de la tabla.
% ⚠️ Si alguna fila cambia, cambiarla tambien en doc/SoTA/comparative_analysis.md.
\begin{table}[t]
    \centering
    \small
    \begin{tabular}{@{}lcccc@{}}
        \toprule
        \textbf{Trabajo} & \textbf{Pre-ejec.} & \textbf{Magnitud} & \textbf{Multi-salida} & \textbf{Deriva} \\
        \midrule
        ZNE, PEC, M3                            & no & sí & no & no \\
        \citeNP{liao2024machine}                & no & sí & no & parcial \\
        GTraQEM, QEMFormer                      & no & sí & sí & no \\
        mapomatic, \citeNP{hartnett2024learning}& sí & no & no & no \\
        \citeNP{wang2025quest}                  & sí & sí & no & no \\
        \citeNP{liu2026output}                  & sí & sí & no & no \\
        \citeNP{aktar2024graph}                 & sí & sí & no & no \\
        \citeNP{martyniuk2025benchmarking}      & sí & sí & no & no \\
        \citeNP{tudisco2025graph}               & sí & no & no & no \\
        \midrule
        \textbf{Este trabajo}                   & \textbf{sí} & \textbf{sí} & \textbf{sí} & \textbf{sí} \\
        \bottomrule
    \end{tabular}

    \caption{El estado del arte frente a los cuatro rasgos que definen esta propuesta.}
    \label{tab:hueco}

    {\footnotesize Fuente: elaboración propia.}
\end{table}

El recorrido del estado del arte deja ver que las investigaciones están partidas en dos mitades. Una
corrige bien, pero solo \textbf{después} de haber gastado la máquina. La otra decide
\textbf{antes}, pero lo que predice sirve para descartar un circuito o para elegir dónde
ejecutarlo, no para corregir su resultado.

La Tabla~\ref{tab:hueco} contrasta el corpus con los cuatro rasgos que definen esta propuesta, y
\textbf{ninguna fila reúne los cuatro}:

\begin{enumerate}
    \item \textbf{Pre-ejecución}: la entrada del modelo nunca incluye el resultado medido. Es la
          frontera entre las dos mitades, y la que decide si el método ahorra tiempo de
          procesador o lo consume.
    \item \textbf{Magnitud}: predecir cuánto se degrada, y no solo si conviene ejecutar.
    \item \textbf{Multi-salida} sobre varios observables a la vez, en lugar de un escalar único.
          Los trabajos pre-ejecución revisados predicen todos una sola cantidad.
    \item \textbf{El paso del tiempo como eje}, con días de calibración reales y disjuntos entre
          entrenamiento y prueba. Hay evidencia publicada de que el hardware deriva a saltos, y
          aun así ningún trabajo de la literatura lo incluye en su evaluación.
\end{enumerate}
